\section{Experimental protocol}
\subsection{Historical and modern separation}
Historical results come from the legacy freeze: original AC--MOT, V1 Trial 24,
V2 Trial 22, historical VisDrone test-dev, and historical UAVDT transfer.
Modern results come from tag \texttt{acmot-v3-exp-2026-10-04-freeze} and its
archived experiment package. The protocols are not pooled because their
detectors, operating grids, evaluation details, and held-out roles differ.

\subsection{Modern detector, splits, and transfer}
The modern detector is \texttt{dronefreak/visdrone-yolo11m} trained on
VisDrone2019-DET at 640 pixels; the checkpoint SHA256 is \checkpointsha.
Development uses eight VisDrone2019-MOT calibration sequences and four
sequence-level folds. The modern VisDrone read uses seven sequences and 71,830
ground-truth boxes. It is the project's class-agnostic protocol, not an
official VisDrone leaderboard evaluation, and it is disclosed as a second
validation read because a static control had already read the sequences.

UAVDT transfer uses 20 sequences, 16,592 frames, and 49,776 detector forward
passes from three action-specific cache pairs. The frozen policy is applied
without recalibration or retuning, once per tracker host, so no bootstrap or
cross-validation claim is attached.

\subsection{Trackers, baseline, and metrics}
ByteTrack uses the frozen archive defaults. OATrack is the repository's frozen
reimplementation with tracker-side \texttt{min\_conf=0.40}; this is separate
from the detector-side v3 confidence floor. The declared modern baseline is
\texttt{r1536\_n70}: resolution 1536, NMS IoU 0.70, and no detector filter
above the cache floor 0.01. V3 uses the frozen action map and thresholds; no
detector, tracker, or operating choice changes after the freeze.

We report HOTA, MOTA, IDF1, IDS, FP, FN, precision, and recall. HOTA balances
detection and association, MOTA exposes the combined FP/FN/IDS effect, and
IDF1 measures identity precision and recall. Baseline and v3 are paired by
sequence. The archived modern bootstrap uses 5000 paired resamples with seed
42; it is not applied to the single-run UAVDT transfer.

\subsection{Runtime and evaluation handling}
V3 runtime is measured on a Tesla T4 over 200 frames with causal scene
analysis, detector execution at the selected resolution, and tracker update.
Mean latency, P95 latency, FPS, detector time, tracker time, and controller
overhead are recorded. The baseline harness does not include image-read work
identically, so the comparison is reported with that caveat and is not called a
perfectly matched speedup. The analysis does not claim a 30-FPS deployment
target.

Detector outputs are cached at the required resolution/NMS combinations for
reproducibility. The project maps classes to its documented class-agnostic
vehicle/tracking protocol before TrackEval-style scoring; these values must
therefore not be described as official leaderboard metrics. The archive records
the detector hash, controller thresholds, action map, tracker descriptions,
calibration folds, forward-pass count, and source checksums. It also records
that frozen Trial 7 is not the literal argmax of the documented selection rule
under either eligibility interpretation; the paper reports this deviation and
does not rerun the search.
